% !TeX root = ../../Book3.tex
% 纲要标题：### F.2 Mora 弱正规形的动机

\section{Mora 弱正规形的动机}
\label{b3:sec:F02}

局部序中的首项消去可能产生无穷过程，有限输入并不自动保证普通除法终止。在 \(k[x]_{(x)}\) 中，\(g=x-x^2\) 确实生成 \((x)\)，\(\{g\}\) 已是标准基，\(x\in(g)\) 也有有限证书；只按首项反复约化 \(x\) 却仍会产生无穷过程。Mora 方法允许在被除式前乘一个局部单位，并用次数差控制新的约化步骤。

% 纲要细目（与计划纲要同步）：
% * 用 \(g=x-x^2\) 区别正确的成员关系、正确标准基与朴素约化不终止三件事
% * 单位乘子证书与完整 Mora 弱正规形的首项界分开；临时约化元须由原生成元表示核验
% * 只展示局部次数序的机制与有限核验；一般算法正确性、终止性和更广混合序策略准确定位原文，保留专题范围
% #### F.2.1 局部约化中的无穷过程
% #### F.2.2 单位乘子、弱正规形证书与理想成员核验
% #### F.2.3 Mora 方法的算例与进一步讨论

\subsection{局部约化中的无穷过程}
\label{b3:subsec:F02:01}

在 \(k[x]\) 上取局部序 \(1>x>x^2>\cdots\)，令 \(g=x-x^2\)。若照搬“用首项 \(x\) 消去首项”的规则，便有
\[
x-g=x^2,\qquad x^2-xg=x^3,\qquad
x^r-x^{r-1}g=x^{r+1}.
\]
余项不断向高次移动，却永远不为零。每一步都合法，但“次数越来越高”不意味着作为多项式的过程已经结束。

\ProseParagraphBreak
在形式幂级数环中，这个过程对应
\[
x=(1+x+x^2+\cdots)g.
\]
这个等式按 \(x\)-进拓扑有意义，但商系数不是多项式。在局部有理函数环中，同一个结论有一个有限表示
\[
(1-x)x=g,\qquad 1-x\in k[x]_{(x)}^\times.
\]
局部计算因此有一个自然目标：允许把输入乘以一个可识别的单位，以换取有限的多项式等式。它既保留局部成员关系，又避免将无限级数展开到底。

\begin{remark}
上述不终止并不说明 \(g\) 不是标准基。事实上它在局部环中生成 \((x)\)，首项也是 \(x\)，故已经是标准基。问题出在求余过程，而不是标准基缺少一个首项。这正是局部理论必须另设正规形方法的原因。
\end{remark}

\subsection{单位乘子、弱正规形证书与理想成员核验}
\label{b3:subsec:F02:02}

以下“弱正规形证书”只保留成员核验所需的条件：有限多项式等式、可逆乘子及余项首单项式的不可约性。Mora 弱正规形还附带满足首项上界的标准表示；求出它的算法则须规定约化元选择及暂存规则。这里先定义足以核验成员关系的简化证书。

\begin{definition}\label{b3:def:F-weak-normal-form}
设 \(k\) 为域，\(n\ge1\)，\(S=k[x_1,\ldots,x_n]\)，固定单项式次序 \(>\)，令 \(U_{>}=\{u\in S:\operatorname{LM}_{>}(u)=1\}\)、\(R=U_{>}^{-1}S\)。给定由非零多项式组成的有限列 \(G=(g_1,\ldots,g_s)\) 和 \(f\in S\)，若 \(u\in U_{>}\)、\(a_1,\ldots,a_s,h\in S\) 满足
\[
uf=\sum_{i=1}^s a_ig_i+h,
\]
且 \(h=0\)，或 \(\operatorname{LM}_{>}(h)\) 不被任何
\(\operatorname{LM}_{>}(g_i)\) 整除，则称 \((u,a_1,\ldots,a_s,h)\) 及这份等式为 \(f\) 关于 \(G\) 的一份\term{弱正规形证书}[weak normal form certificate]。此处只要求余项的首单项式不可约，不要求其全部项都不可约。
\end{definition}

标准表示还包含首项上界：当 \(f\ne0\) 时，要求每个非零 \(a_ig_i\) 满足
\[
\operatorname{LM}_{>}(a_ig_i)
\le_{>}\operatorname{LM}_{>}(f).
\]
因为 \(\operatorname{LM}_{>}(u)=1\)，等式左边与 \(f\) 有相同的首单项式；这一上界保证生成元组合没有先引入更大的首项再相消。若只有多项式等式，表示中的各项可以带有比目标更大的首项，最后依靠相消才得到目标；成员核验仍然有效，却不能据此控制临界对的首项抵消。Mora 约化同时追踪正规形与标准表示，后者是临界对判据所需的信息，见\cite[Definitions 5--6, Section 3.1]{MoraPfisterTraverso1989TangentCone}及\cite[Section 2, Theorem 2.14]{Aschenbrenner2005ReductionStandardBases}。上面所定义的证书只保留成员核验条件，并不把任意有限表示都视为满足首项上界的 Mora 弱正规形。

\begin{proposition}\label{b3:prop:F-weak-membership}
设 \(k\) 为域，\(n\ge1\)，\(S=k[x_1,\ldots,x_n]\)，\(>\) 为乘法相容的单项式全序，\(R=S_{>}\)。设由非零多项式组成的有限列 \(G=(g_i)\) 是理想 \(I\subset R\) 的标准基，并给定 \(f\in S\) 的弱正规形证书 \(uf=\sum_i a_i g_i+h\)，其中 \(u\in U_{>}\)、\(a_i,h\in S\)。则 \(f\in I\) 当且仅当 \(h=0\)。
\end{proposition}
\begin{proof}
若 \(h=0\)，由于 \(u\) 在 \(R\) 中可逆，立即得到 \(f\in I\)。若 \(f\in I\)，则 \(h=uf-\sum_i a_ig_i\in I\)。若 \(h\ne0\)，标准基条件迫使其首单项式被某个 \(\operatorname{LM}_{>}(g_i)\) 整除，与证书对余项的要求矛盾。
\end{proof}

给定这份证书，只需核验多项式等式、单位条件和余项首单项式的不可约性，就能判定成员关系。证书不必唯一；求出证书的算法还需另外规定约化元选择及终止机制。在原点局部环中，单位条件就是 \(u(0)\ne0\)。

\ProseParagraphBreak
例如对 \(g=x-x^2\)，\(f=x\) 的证书就是 \(u=1-x,a_1=1,h=0\)。相反，对 \(f=1\) 可取 \(u=1,a_1=0,h=1\)；首单项式 \(1\) 不被 \(x\) 整除，证书中的非零余项说明 \(1\notin(g)k[x]_{(x)}\)。

\subsection{Mora 方法的算例与进一步讨论}
\label{b3:subsec:F02:03}

这里通过一个有限算例解释暂存机制；域上的多项式输入及局部次数序下，完整 Mora 约化的选择规则、正确性与终止证明采用本小节末所引文献的结果。Mora 方法的一个关键想法是把某些中间余项暂时加入可用的约化集合，避免永远重复最初的约化。选择约化元时可参考
\[
\operatorname{ecart}(f)=\deg(f)-\deg(\operatorname{LM}_{>}(f)),
\qquad f\ne0.
\]
在局部次数序中，它衡量多项式最高次数与最低出现次数之间的差。若只选择首项可整除而忽略后面的高次项，正可能重复产生更高次余项。
\symidx{\(\operatorname{ecart}(f)\)：多项式的次数与首单项式次数之差}

\begin{example}\label{b3:exm:F-mora-one-step}
继续取 \(f=x\)、\(G=\{x-x^2\}\)。起初
\[
\operatorname{ecart}(x)=0,\qquad
\operatorname{ecart}(x-x^2)=1.
\]
使用差值更大的 \(x-x^2\) 前，先把当前余项 \(x\) 暂存。第一次相减得到 \(h=x^2\)，接着可使用暂存的 \(x\) 消去它，得到 \(h=0\)。暂存只扩充可用的约化集合，原理想仍由 \(G\) 生成；这一动作本身不证明 \(x\in(g)\)，也不能直接把 \(x\) 当作新取得的理想生成元。必须把两步还原成关于原输入和原生成元的等式：
\[
h_1=f-g=x^2,\qquad h_2=h_1-xf=0.
\]
合并得 \((1-x)f=g\)，正好产生一个单位乘子。因 \(1-x\) 的常数项为 \(1\)，这份有限证书有效。
\end{example}

这个算例说明暂存余项如何转化为单位乘子，也说明必须追踪原输入和原生成元。一般 Mora 算法须精确规定约化元选择、暂存规则和首项上界；在域上的多项式输入及局部次数序下，完整终止证明结合次数差的控制、Dickson 引理与同次单项式的有限性，见\cite[Section 3.1, Lemmas 10--11 and Corollary 2]{MoraPfisterTraverso1989TangentCone}。本节核验了成员证书及上述两步算例，一般正确性与终止性采用该文献的证明。单凭暂存余项不能推出任意输入上的终止性；这里基于标量次数差与同次单项式有限性的说明只涉及局部次数序，对非次数相容的局部序及混合序，还须按具体次序规定约化策略并核验相应的终止论证，参见同一文献的第3.2节。

\ProseParagraphBreak
形式除法采取另一种完成方式：不追求有限步停止，而证明每个固定次数的系数最终稳定。两种方法可以计算相同局部对象，但一个以单位乘子给出有限表示，一个以拓扑保证无限和有意义。后者得到的低次关系还能组成初始齐次理想，用于计算切锥。
